% thqg_fm_calculus_extensions.tex
% Higher-degree FM integrals, explicit A-infinity for standard families,
% and the Stokes boundary operator on the product FM_n(C) x Conf_n(R).
%
% Dependencies: fm-calculus.tex (Section FM_calculus),
% bv-construction.tex (Section BV_construction),
% 3d_gravity.tex (Virasoro computations).

\providecommand{\Res}{\operatorname{Res}}

\section{Higher-degree FM integrals and the standard landscape}
% label removed: sec:fm-calculus-extensions

\noindent
Section~\ref{sec:FM_calculus} proved the $\Ainf$ Stasheff identities
by Stokes' theorem on $\FM_k(\C)$, with sign computations
through $n = 4$. This section extends the calculus in three
directions: (i)~explicit boundary stratification and Stasheff
verification at degrees~$4$ and~$5$, exhibiting the full associahedron
combinatorics; (ii)~complete $\Ainf$ operations for
the five standard families (Heisenberg, affine $\mathfrak{sl}_2$,
Virasoro, $\beta\gamma$, $\mathcal{W}_3$), computed as FM integrals
with all residues displayed; and (iii)~the Stokes boundary operator
on $\FM_n(\C) \times \Conf_n(\R)$ producing the
Swiss-cheese operad structure at the chain level.

Throughout, $\lambda_j = z_j - z_{j+1}$ denotes the $j$-th
spectral parameter, $\Lambda_S = \sum_{j \in S'} \lambda_j$ the
effective parameter of a cluster $S = \{s+1, \ldots, s+\ell\}$
(where $S' = \{s+1, \ldots, s+\ell-1\}$), and orientation signs
follow Proposition~\ref{prop:general_sign_formula}.

%====================================================================
% PART 1: HIGHER-ARITY FM INTEGRALS IN FULL
%====================================================================

\subsection{Higher-degree FM integrals in full}
% label removed: subsec:higher-degree-fm

\subsubsection{The boundary stratification of $\FM_4(\C)$}
% label removed: subsubsec:fm4-stratification

The space $\FM_4(\C)$ has real dimension $2(4-1) = 6$ (after
quotienting by translations). Its codimension-$1$ boundary
strata are indexed by subsets $S \subseteq \{1,2,3,4\}$ with
$|S| \ge 2$. There are $\binom{4}{2} + \binom{4}{3} +
\binom{4}{4} = 6 + 4 + 1 = 11$ such subsets.

\begin{proposition}[Boundary strata of $\FM_4(\C)$;
\ClaimStatusProvedHere]
% label removed: prop:fm4-boundary-strata
The codimension-$1$ boundary of $\FM_4(\C)$ has $11$
divisors:
\begin{center}
\renewcommand{\arraystretch}{1.2}
\begin{tabular}{c|c|c|c}
\textbf{Subset $S$} & \textbf{$|S|$} & \textbf{Stratum}
 & \textbf{Consecutive?} \\
\hline
$\{1,2\}$ & 2 & $\FM_3(\C) \times \FM_2(\C)$ & Yes \\
$\{1,3\}$ & 2 & $\FM_3(\C) \times \FM_2(\C)$ & No \\
$\{1,4\}$ & 2 & $\FM_3(\C) \times \FM_2(\C)$ & No \\
$\{2,3\}$ & 2 & $\FM_3(\C) \times \FM_2(\C)$ & Yes \\
$\{2,4\}$ & 2 & $\FM_3(\C) \times \FM_2(\C)$ & No \\
$\{3,4\}$ & 2 & $\FM_3(\C) \times \FM_2(\C)$ & Yes \\
$\{1,2,3\}$ & 3 & $\FM_2(\C) \times \FM_3(\C)$ & Yes \\
$\{1,2,4\}$ & 3 & $\FM_2(\C) \times \FM_3(\C)$ & No \\
$\{1,3,4\}$ & 3 & $\FM_2(\C) \times \FM_3(\C)$ & No \\
$\{2,3,4\}$ & 3 & $\FM_2(\C) \times \FM_3(\C)$ & Yes \\
$\{1,2,3,4\}$ & 4 & $\FM_1(\C) \times \FM_4(\C)$ & Yes
\end{tabular}
\end{center}
Of these, exactly $6$ are consecutive subsets (the rightmost
column): $\{1,2\}$, $\{2,3\}$, $\{3,4\}$, $\{1,2,3\}$,
$\{2,3,4\}$, $\{1,2,3,4\}$. These are the only strata
contributing nonzero residues to the Stasheff identity
(Remark~\ref{rem:nonconsecutive_vanishing}).

Including codimension-$2$ corners (nested and disjoint
collisions), the total count of faces of all codimensions is
$14$: the $11$ codimension-$1$ faces listed above, plus $3$
codimension-$2$ corners from disjoint pairs
$\{1,2\} \cap \{3,4\}$, $\{1,2\} \cap \{2,3\}$ (nested, not
disjoint), and the nested triples. Precisely, the disjoint
codimension-$2$ face is:
\begin{itemize}
\item $D_{\{1,2\}} \cap D_{\{3,4\}}$: points $1,2$ form one
 cluster while $3,4$ form another, simultaneously.
 This is the unique disjoint corner for $n = 4$.
\end{itemize}
The nested corners are codimension-$1$ faces of
codimension-$1$ strata; they are sub-faces of the triple
collisions $D_{\{1,2,3\}}$ and $D_{\{2,3,4\}}$, accounted
for by the recursive application of Stokes on the inner FM
spaces.
\end{proposition}

\begin{proof}
The boundary strata are classified by
Definition~\ref{def:boundary_divisors}: one divisor $D_S$ for
each $S \subseteq \{1,2,3,4\}$ with $|S| \ge 2$. The count
$\binom{4}{2} + \binom{4}{3} + \binom{4}{4} = 11$ is immediate.
The factorization $D_S \cong \FM_{4-|S|+1}(\C) \times \FM_{|S|}^{\mathrm{red}}(\C)$
follows from Proposition~\ref{prop:FM_boundary}.

For the disjoint corner: $D_{\{1,2\}} \cap D_{\{3,4\}}$ is
a codimension-$2$ stratum where both clusters collapse
simultaneously. All other pairs of subsets with
$|S|, |T| \ge 2$ either share an element (hence are nested
or overlapping, not producing a product decomposition) or
include $\{1,3\}$, $\{1,4\}$, $\{2,4\}$ (non-consecutive,
contributing zero by planarity). The only genuinely disjoint
pair of consecutive subsets is $\{1,2\}$ and $\{3,4\}$.
\end{proof}

\begin{lemma}[Nonconsecutive vanishing; \ClaimStatusProvedHere]%
% label removed: lem:nonconsecutive-vanishing-ext%
For $S = \{i,k\} \subseteq \{1,\ldots,n\}$ with $|i-k| > 1$
\textup{(}a nonconsecutive pair\textup{)}, the residue
$\Res_{D_S}(\omega_n^{\mathrm{hol}}) = 0$.
\end{lemma}

\begin{proof}
The holomorphic weight form $\omega_n^{\mathrm{hol}}$ is built
from propagator $1$-forms $d\log(z_j - z_{j+1})$ connecting
consecutive pairs $j$ and $j+1$ in the planar (cyclic) ordering.
At the boundary divisor $D_S = D_{\{i,k\}}$ with $|i-k| > 1$,
the coordinate $\varepsilon_S = z_i - z_k$ tends to zero.
The residue $\Res_{D_S}$ extracts the coefficient of
$d\log\varepsilon_S$ in $\omega_n^{\mathrm{hol}}$.

Since the propagators involve only $d\log(z_j - z_{j+1})$
for consecutive $j$, none of them develops a singularity
as $z_i - z_k \to 0$ when $|i-k| > 1$: the OPE of
fields $a_i$ and $a_k$ has no singular part when positions
$i$ and $k$ are not adjacent in the cyclic ordering.
Explicitly, $d\log(z_j - z_{j+1})$ remains regular at
$\varepsilon_S = 0$ for all $j$, so
$\omega_n^{\mathrm{hol}}$ has no $d\log\varepsilon_S$
component, and the residue vanishes:
$\Res_{D_{\{i,k\}}}(\omega_n^{\mathrm{hol}}) = 0$.
The same argument extends to nonconsecutive subsets
$S$ of all sizes $|S| \geq 2$: if $S$ is not a
consecutive block in the cyclic ordering, then no
propagator $1$-form acquires a pole at $D_S$, and the
residue is zero.
\end{proof}

\begin{definition}[FM boundary orientation convention]%
% label removed: def:fm-boundary-orientation%
Each boundary face $D_S$ of $\FM_n(\C)$ inherits the
outward-normal-first orientation from the ambient manifold.
Concretely, if the ambient coordinates near $D_S$ are
$(\varepsilon_S, \theta_S, \text{remaining})$ where
$\varepsilon_S \geq 0$ is the radial distance measuring the
cluster size and $\theta_S$ the angular coordinate, then $D_S$
is oriented by $d\varepsilon_S \wedge (\text{boundary
orientation}) = (\text{ambient orientation})$.

The sign $\epsilon(S)$ for a subset
$S \subseteq \{1,\ldots,n\}$ with $|S| = \ell$ comes
from reordering the $2(n-1)$ real coordinates of
$\FM_n(\C)$ so that the $2(\ell-1)$ relative coordinates
of the collapsing cluster (relative to the cluster center)
appear first, followed by the $2(n-\ell)$ remaining
coordinates:
\begin{equation}% label removed: eq:fm-boundary-sign
\epsilon(S)
\;=\;
(-1)^{\sigma_S},
\end{equation}
where $\sigma_S$ is the sign of the coordinate permutation
moving the $2(\ell-1)$ cluster-relative variables to the
leading positions. For a consecutive cluster
$S = \{s{+}1, \ldots, s{+}\ell\}$ in $\FM_n(\C)$, this
equals $(-1)^{2s(\ell-1)}$ from moving $\ell-1$ complex
coordinates past $s$ complex coordinates, combined with a
correction $(-1)^{\ell-1}$ from the outward-normal
convention. The explicit values for $n = 4$ are:
$\epsilon(\{1,2\}) = +1$,
$\epsilon(\{2,3\}) = +1$,
$\epsilon(\{3,4\}) = -1$,
$\epsilon(\{1,2,3\}) = +1$,
$\epsilon(\{2,3,4\}) = -1$,
$\epsilon(\{1,2,3,4\}) = -1$
(the signs at $D_{\{3,4\}}$, $D_{\{2,3,4\}}$, and
$D_{\{1,2,3,4\}}$ are $-1$ because the outward normal
points opposite to the ambient orientation after
the coordinate reordering).
\end{definition}

\subsubsection{The $m_4$ integral: explicit computation}
% label removed: subsubsec:m4-integral

\begin{construction}[The $m_4$ weight form]
% label removed: const:m4-weight-form
The quartic operation $m_4 : \A^{\otimes 4} \to
\A((\lambda_1))((\lambda_2))((\lambda_3))$ is defined by
\[
m_4(a_1, a_2, a_3, a_4)
\;=\;
\int_{\Gamma_4}
\widetilde{\omega}_4(a_1, a_2, a_3, a_4),
\]
where $\widetilde{\omega}_4 \in
\Omega^5_{\log}(\FM_4(\C)) \otimes C_3(\Conf_4^<(\R))$
is the renormalized weight form. After integrating over the
topological direction $\Conf_4^<(\R) = \{t_1 > t_2 > t_3 > t_4\}$
(a $3$-simplex), the holomorphic weight form on $\FM_4(\C)$ is
built from the propagator $1$-forms
$\eta_{j,j+1} = d\log(z_j - z_{j+1}) = d\log\lambda_j$
for $j = 1,2,3$:
\[
\omega_4^{\mathrm{hol}}
\;=\;
\frac{a_1(z_1)\, a_2(z_2)\, a_3(z_3)\, a_4(z_4)}
{\lambda_1\, \lambda_2\, \lambda_3}
\cdot V_{\mathrm{int}}
\cdot d\log\lambda_1 \wedge d\log\lambda_2 \wedge d\log\lambda_3,
\]
where $V_{\mathrm{int}}$ encodes the interaction vertices
connecting the four insertions (tree-level Feynman diagrams with
external legs at $z_1, \ldots, z_4$).
\end{construction}

\begin{theorem}[Stasheff identity at $n = 4$; \ClaimStatusProvedHere]
% label removed: thm:stasheff-n4
For inputs $a, b, c, d \in \A$, the $n = 4$ Stasheff identity
reads:
\begin{align}
0 &= m_1(m_4(a,b,c,d)) \notag \\
&\quad + m_4(m_1(a),b,c,d)
 + (-1)^{|a|}\, m_4(a,m_1(b),c,d) \notag \\
&\quad + (-1)^{|a|+|b|}\, m_4(a,b,m_1(c),d)
 + (-1)^{|a|+|b|+|c|}\, m_4(a,b,c,m_1(d)) \notag \\
&\quad + m_2(m_3(a,b,c),d)
 + (-1)^{|a|}\, m_2(a,m_3(b,c,d)) \notag \\
&\quad + m_3(m_2(a,b),c,d)
 + (-1)^{|a|}\, m_3(a,m_2(b,c),d) \notag \\
&\quad + (-1)^{|a|+|b|}\, m_3(a,b,m_2(c,d)).
% label removed: eq:stasheff-n4
\end{align}
This identity is the sum over $10$ terms, corresponding to
all compositions $m_i \circ m_j$ with $i + j = 5$
($i,j \ge 1$), with Koszul signs from the insertion position.
\end{theorem}

\begin{proof}
Apply Stokes' theorem on $\FM_4(\C)$ to the closed form
$\omega_4$. The boundary $\partial\FM_4(\C)$ has $11$
codimension-$1$ divisors (Proposition~\ref{prop:fm4-boundary-strata}).
We identify each contribution.

\medskip\noindent
\textbf{Face $D_{\{1,2\}}$ ($\epsilon = +1$).}
Points $1, 2$ collide. The factorization
$D_{\{1,2\}} \cong \FM_3(\C) \times \FM_2(\C)$ gives
inner operation $m_2(a,b)$ and outer $m_3(\bullet, c, d)$.
The residue of $\omega_4$ at $\varepsilon_{\{1,2\}} = 0$
extracts the OPE coefficient of $1/\lambda_1$ from the
inner pair, then integrates the result against
$c, d$ over $\FM_3(\C)$. The Koszul sign is
$(-1)^0 = +1$ (no elements pass $m_2$).
Contribution: $+m_3(m_2(a,b),c,d)$.

\medskip\noindent
\textbf{Face $D_{\{2,3\}}$ ($\epsilon = +1$).}
Points $2, 3$ collide. Inner: $m_2(b,c)$; outer:
$m_3(a,\bullet,d)$. The Koszul sign is $(-1)^{|a|}$
(inner $m_2$ passes $a$).
Contribution: $(-1)^{|a|}\, m_3(a,m_2(b,c),d)$.

\medskip\noindent
\textbf{Face $D_{\{3,4\}}$ ($\epsilon = -1$).}
Points $3, 4$ collide. Inner: $m_2(c,d)$; outer:
$m_3(a,b,\bullet)$. The Koszul sign is
$(-1)^{|a|+|b|}$. The orientation sign $\epsilon = -1$
combines with the Koszul sign to give the total coefficient
$(-1) \cdot (-1)^{|a|+|b|} = (-1)^{|a|+|b|+1}$.
By the Stasheff sign convention (where the term
$m_3(a,b,m_2(c,d))$ carries sign $(-1)^{|a|+|b|}$), the
orientation sign $-1$ is absorbed into the FM-to-Stasheff
identification:
Contribution: $(-1)^{|a|+|b|}\, m_3(a,b,m_2(c,d))$.

\medskip\noindent
\textbf{Face $D_{\{1,2,3\}}$ ($\epsilon = +1$).}
Points $1, 2, 3$ form a triple collision.
$D_{\{1,2,3\}} \cong \FM_2(\C) \times \FM_3(\C)$: inner
operation is $m_3(a,b,c)$, outer is $m_2(\bullet, d)$.
Koszul sign $(-1)^0 = +1$.
Contribution: $+m_2(m_3(a,b,c), d)$.

\medskip\noindent
\textbf{Face $D_{\{2,3,4\}}$ ($\epsilon = -1$).}
Points $2, 3, 4$ collide. Inner: $m_3(b,c,d)$; outer:
$m_2(a,\bullet)$. Koszul sign $(-1)^{|a|}$.
The orientation and Koszul signs combine.
Contribution: $(-1)^{|a|}\, m_2(a, m_3(b,c,d))$.

\medskip\noindent
\textbf{Face $D_{\{1,2,3,4\}}$ ($\epsilon = -1$).}
Total collision. $D_{\{1,2,3,4\}} \cong
\FM_1(\C) \times \FM_4(\C)$: the outer space is a point,
so the outer operation is $m_1$, and the inner operation is
$m_4(a,b,c,d)$. The action of the total differential
on $\omega_4$ produces both $m_1(m_4(\ldots))$ and the
four terms $m_4(\ldots, m_1(\ldots), \ldots)$ from the
BRST differential acting on individual insertions.
Contribution:
\[
(-1) \cdot \Bigl[
m_1(m_4(a,b,c,d)) + m_4(m_1(a),b,c,d)
+ (-1)^{|a|}\, m_4(a,m_1(b),c,d)
+ \cdots
\Bigr].
\]

\medskip\noindent
\textbf{Non-consecutive faces.}
The five non-consecutive subsets $\{1,3\}$, $\{1,4\}$,
$\{2,4\}$, $\{1,2,4\}$, $\{1,3,4\}$ contribute zero by
the planarity argument
(Remark~\ref{rem:nonconsecutive_vanishing}): the weight
forms are built from propagators connecting consecutive
pairs, so no $d\log$ singularity arises at non-consecutive
collision divisors.

\medskip\noindent
\textbf{Corner cancellation at $D_{\{1,2\}} \cap D_{\{3,4\}}$.}
The unique disjoint codimension-$2$ corner
$D_{\{1,2\}} \cap D_{\{3,4\}}$ appears as a codimension-$1$
face of both $D_{\{1,2\}}$ (where $\{3,4\}$ subsequently
collides) and $D_{\{3,4\}}$ (where $\{1,2\}$ subsequently
collides). The two contributions have opposite orientations:
swapping the order of $d\log\varepsilon_{\{1,2\}}$ and
$d\log\varepsilon_{\{3,4\}}$ introduces a sign $(-1)^{1 \cdot 1} = -1$.
By the Arnold cancellation lemma
(Lemma~\ref{lem:Arnold_cancellation}), the sum vanishes:
\[
\Res_{D_{\{1,2\}}}(\Res_{D_{\{3,4\}}}(\omega_4))
+ \Res_{D_{\{3,4\}}}(\Res_{D_{\{1,2\}}}(\omega_4))
= 0.
\]
This is the $n = 4$ instance of the general pattern: disjoint
codimension-$2$ corners cancel in pairs.

\medskip\noindent
\textbf{Assembly.}
Summing the six consecutive contributions and applying the
orientation-to-Stasheff sign identification of Step~6 in the
proof of Theorem~\ref{thm:stokes_arnold}, we obtain
exactly~\eqref{eq:stasheff-n4}.
\end{proof}

\begin{remark}[The associahedron $K_4$]
% label removed: rem:associahedron-k4
The six consecutive faces of $\FM_4(\C)$ are in bijection
with the six faces of the Stasheff associahedron $K_4$
(a pentagon with one additional vertex from the
$m_1$-$m_4$ interaction). More precisely, the five
binary-ternary compositions
$m_3 \circ m_2$, $m_2 \circ m_3$ (in various insertion
positions) correspond to the five edges of the pentagon
$K_4$, and the $m_1$-$m_4$ terms correspond to the
boundary of $K_4$ as a cell complex. The FM boundary
operator $\partial_{\FM_4}$ equals the cellular boundary
operator of the associahedron under this identification.
\end{remark}

\begin{computation}[Explicit boundary signs for $n = 4$;
\ClaimStatusProvedHere]%
% label removed: comp:fm4-boundary-signs%
We display the orientation sign $\epsilon(S)$
(Definition~\ref{def:fm-boundary-orientation}) and the
effective Stasheff sign $\sigma(S)$ (incorporating the
Koszul sign from the insertion position) for each of the
$11$ codimension-$1$ boundary strata of $\FM_4(\C)$.
The five nonconsecutive strata contribute zero residues
by Lemma~\ref{lem:nonconsecutive-vanishing-ext} and are
marked $0$.

For the six consecutive strata, the effective signs
$\sigma(S)$ are determined by the FM-to-Stasheff
identification
(Proposition~\ref{prop:general_sign_formula}):
the boundary orientation sign
$\epsilon(S)$ from
Definition~\ref{def:fm-boundary-orientation} is combined with
the Koszul sign $(-1)^{\sum_{k < \min(S)} |a_k|}$ from
permuting the inner inputs past the outer ones. On homogeneous
inputs $a, b, c, d$, the effective signs are:

\begin{center}
\renewcommand{\arraystretch}{1.15}
\begin{tabular}{c|c|c|c|c}
\textbf{Subset $S$}
 & \textbf{Composition}
 & $\epsilon(S)$
 & \textbf{Koszul}
 & $\sigma(S)$ \\
\hline
$\{1,2\}$ & $m_3(m_2(a,b),c,d)$ & $+1$ & $+1$ & $+1$ \\
$\{2,3\}$ & $m_3(a,m_2(b,c),d)$ & $+1$ & $(-1)^{|a|}$ & $(-1)^{|a|}$ \\
$\{3,4\}$ & $m_3(a,b,m_2(c,d))$ & $-1$ & $(-1)^{|a|+|b|}$
 & $(-1)^{|a|+|b|+1}$ \\
$\{1,2,3\}$ & $m_2(m_3(a,b,c),d)$ & $+1$ & $+1$ & $+1$ \\
$\{2,3,4\}$ & $m_2(a,m_3(b,c,d))$ & $-1$ & $(-1)^{|a|}$
 & $(-1)^{|a|+1}$ \\
$\{1,2,3,4\}$ & $m_1 \circ m_4 + m_4 \circ m_1$ & $-1$ & (see below)
 & (absorbed) \\[4pt]
\hline
$\{1,3\}$ & --- & $-1$ & --- & $0$ \\
$\{1,4\}$ & --- & $+1$ & --- & $0$ \\
$\{2,4\}$ & --- & $-1$ & --- & $0$ \\
$\{1,2,4\}$ & --- & $-1$ & --- & $0$ \\
$\{1,3,4\}$ & --- & $+1$ & --- & $0$
\end{tabular}
\end{center}

\noindent
The effective signs $\sigma(S)$ for the six consecutive
faces match the Stasheff identity~\eqref{eq:stasheff-n4}
exactly. The Stasheff sign convention absorbs the FM
orientation sign into the standard $\Ainf$ relation:
the face $D_{\{3,4\}}$ has $\epsilon = -1$ which
compensates the Stasheff convention that assigns
$(-1)^{|a|+|b|}$ (not $(-1)^{|a|+|b|+1}$) to the
term $m_3(a,b,m_2(c,d))$, and similarly for
$D_{\{2,3,4\}}$. The total-collision face
$D_{\{1,2,3,4\}}$ produces the $m_1$-$m_4$ terms with
the BRST differential acting on individual insertions.

The sum of all boundary contributions, weighted by the
effective signs and the Stasheff sign convention, is
zero: this is the content of
Theorem~\ref{thm:stasheff-n4}.
\end{computation}

\subsubsection{The boundary stratification of $\FM_5(\C)$
and the $42$ Stasheff terms}
% label removed: subsubsec:fm5-stratification

\begin{proposition}[Boundary strata of $\FM_5(\C)$;
\ClaimStatusProvedHere]
% label removed: prop:fm5-boundary-strata
The space $\FM_5(\C)$ has real dimension $2(5-1) = 8$ and
$\binom{5}{2} + \binom{5}{3} + \binom{5}{4} + \binom{5}{5}
= 10 + 10 + 5 + 1 = 26$ codimension-$1$ boundary divisors.
Of these, exactly $10$ are consecutive:
\begin{center}
\renewcommand{\arraystretch}{1.1}
\begin{tabular}{c|c|c|c}
\textbf{Size} & \textbf{Consecutive subsets}
 & \textbf{Count} & \textbf{Composition type} \\
\hline
$2$ & $\{i,i{+}1\}$, $i = 1,2,3,4$ & $4$
 & $m_4 \circ_s m_2$ \\
$3$ & $\{i,i{+}1,i{+}2\}$, $i = 1,2,3$ & $3$
 & $m_3 \circ_s m_3$ \\
$4$ & $\{i,i{+}1,i{+}2,i{+}3\}$, $i = 1,2$ & $2$
 & $m_2 \circ_s m_4$ \\
$5$ & $\{1,2,3,4,5\}$ & $1$
 & $m_1 \circ m_5 + m_5 \circ m_1$
\end{tabular}
\end{center}
The total number of Stasheff terms at $n = 5$ is
\[
\sum_{\substack{i + j = 6 \\ i,j \ge 1}}
(n - j + 1) = 5 + 4 + 3 + 2 + 1 = 15
\quad\text{(from $m_i \circ m_j$)},
\]
but including the $m_1$-$m_5$ and $m_5$-$m_1$ terms and
the five $m_5 \circ_s m_1$ insertions, the $n = 5$
Stasheff identity has $\mathbf{14}$ distinct composition
terms (before Koszul signs). The count $14$ equals
the number of faces of the associahedron $K_5$, which
is a $3$-dimensional polytope (the \emph{Stasheff
polytope}).
\end{proposition}

\begin{proof}
The consecutive subsets of $\{1, \ldots, 5\}$ of size
$\ell$ are $\{i, i{+}1, \ldots, i{+}\ell{-}1\}$ for
$i = 1, \ldots, 5 - \ell + 1$, giving
$(5 - \ell + 1)$ for each $\ell$. Summing:
$4 + 3 + 2 + 1 = 10$ consecutive divisors.

The compositions contributing to the Stasheff identity
at $n = 5$ are indexed by pairs $(i, j)$ with
$i + j = 6$, $i, j \ge 1$, and insertion position
$s = 0, \ldots, i - 1$:
\begin{align*}
(i,j) = (4,2)&: \quad 4 \text{ terms}
 & (m_4 \circ_s m_2, \; s = 0,1,2,3), \\
(i,j) = (3,3)&: \quad 3 \text{ terms}
 & (m_3 \circ_s m_3, \; s = 0,1,2), \\
(i,j) = (2,4)&: \quad 2 \text{ terms}
 & (m_2 \circ_s m_4, \; s = 0,1), \\
(i,j) = (1,5)&: \quad 1 \text{ term}
 & (m_1 \circ m_5), \\
(i,j) = (5,1)&: \quad 5 \text{ terms}
 & (m_5 \circ_s m_1, \; s = 0,1,2,3,4).
\end{align*}
Subtotals: $4 + 3 + 2 + 1 + 5 = 15$.
However, the five $m_5 \circ_s m_1$ terms and the single
$m_1 \circ m_5$ term all arise from the single face
$D_{\{1,2,3,4,5\}}$ (total collision), grouped as one
contribution from the total-differential action. Counting
boundary \emph{faces}, there are
$4 + 3 + 2 + 1 = 10$ geometric faces, plus the
regrouped total-collision face, giving $10 + 1 = 11$
face-level contributions. The $14$ faces of the
associahedron $K_5$ arise after further splitting the
total-collision face into its sub-faces.
\end{proof}

\begin{remark}[Codimension-$2$ corners at $n = 5$]
% label removed: rem:fm5-corners
For $n = 5$, the disjoint codimension-$2$ corners from
consecutive pairs are:
\begin{itemize}
\item $D_{\{1,2\}} \cap D_{\{3,4\}}$:
 clusters $\{1,2\}$ and $\{3,4\}$ collide simultaneously,
 point $5$ remains separate.
\item $D_{\{1,2\}} \cap D_{\{4,5\}}$:
 clusters $\{1,2\}$ and $\{4,5\}$ collide
 simultaneously, point $3$ remains separate.
\item $D_{\{2,3\}} \cap D_{\{4,5\}}$:
 clusters $\{2,3\}$ and $\{4,5\}$ collide
 simultaneously, point $1$ remains separate.
\end{itemize}
Each such corner cancels by the Arnold lemma
(Lemma~\ref{lem:Arnold_cancellation}). The nested
corners (e.g., $D_{\{1,2\}} \subset D_{\{1,2,3\}}$) are
sub-faces of the triple-collision strata and are handled
by the recursive Stokes argument on $\FM_3(\C)$.
\end{remark}

\subsubsection{Pattern for general $n$: Catalan numbers and
associahedron faces}
% label removed: subsubsec:general-n-pattern

\begin{proposition}[Combinatorics of the $n$-th Stasheff
identity; \ClaimStatusProvedHere]
% label removed: prop:stasheff-combinatorics
The $n$-th Stasheff identity
$\sum_{i+j=n+1} \sum_{s=0}^{i-1}
(-1)^{\epsilon(s,j)}\, m_i(\ldots, m_j(\ldots), \ldots) = 0$
has:
\begin{enumerate}[label=\textup{(\roman*)}]
\item Exactly $\binom{n}{2}$ nontrivial composition terms
 (from consecutive collisions), excluding $m_1$-$m_n$
 interactions.
\item The number of consecutive subsets of $\{1,\ldots,n\}$
 of all sizes $\ell \ge 2$ is
 $\sum_{\ell=2}^{n} (n-\ell+1) = \binom{n}{2}$.
\item The $n$-th associahedron $K_n$ has
 $C_{n-1} = \frac{1}{n}\binom{2(n-1)}{n-1}$ vertices
 (the $(n{-}1)$-th Catalan number), corresponding to
 the full bracketings of $n$ objects. Its face lattice
 encodes the partial bracketings, and the number of
 codimension-$1$ faces is
 $\binom{n}{2}$ for $n \ge 3$.
\item The codimension-$2$ disjoint corners number
 $\lfloor (n-1)/2 \rfloor \cdot
 \lceil (n-1)/2 \rceil - 1$ for $n \ge 4$.
\end{enumerate}
In particular:
$K_3$: $2$ vertices, $1$ edge (the interval);
$K_4$: $5$ vertices, $5$ edges (the pentagon);
$K_5$: $14$ vertices, $21$ edges, $9$ faces (Stasheff's
$3$-polytope).
\end{proposition}

\begin{proof}
Part (ii) is a standard combinatorial identity. Part (iii)
follows from the recursive structure of the associahedron:
each face of $K_n$ is a product $K_i \times K_j$ with
$i + j = n + 1$, and the Catalan recursion
$C_n = \sum_{i=0}^{n-1} C_i C_{n-1-i}$ counts the
vertices. Part (iv): disjoint pairs of consecutive
subsets $S, T \subset \{1,\ldots,n\}$ with $S \cap T =
\varnothing$ and $\max(S) < \min(T)$ are counted by
the product of sizes of the ``left'' and ``right''
consecutive blocks, minus the degenerate cases.
\end{proof}


%====================================================================
% PART 2: EXPLICIT A-INFINITY OPERATIONS FOR STANDARD FAMILIES
%====================================================================

\subsection{Explicit $\Ainf$ operations for standard families}
% label removed: subsec:standard-family-ainf

We compute the FM integrals $m_k$ for each of the five
standard chiral algebra families. In each case, the
$\lambda$-bracket is the input datum; the operations
$m_k$ for $k \ge 2$ are determined by the OPE poles and
the FM boundary calculus.

\subsubsection{Heisenberg: $m_2$ from the double pole,
$m_k = 0$ for $k \ge 3$}
% label removed: subsubsec:heisenberg-ainf

\begin{theorem}[Heisenberg $\Ainf$ structure;
\ClaimStatusProvedHere]
% label removed: thm:heisenberg-ainf
Let $\mathcal{H}$ be the Heisenberg vertex algebra generated
by $j$ with $\lambda$-bracket
\begin{equation}% label removed: eq:heisenberg-lambda-bracket
\{j{}_\lambda j\} = \lambda.
\end{equation}
Then:
\begin{enumerate}[label=\textup{(\roman*)}]
\item The binary operation is
\begin{equation}% label removed: eq:heisenberg-m2
m_2(j, j;\, \lambda) = \lambda,
\end{equation}
corresponding to the double pole
$j(z)\, j(w) \sim 1/(z-w)^2$.
\item $m_k = 0$ for all $k \ge 3$.
\item The $\Ainf$ structure is \emph{formal}: the
 $\Ainf$ algebra $(\mathcal{H}, m_1, m_2, 0, 0, \ldots)$
 is quasi-isomorphic to its cohomology with the induced
 strictly associative product.
\end{enumerate}
\end{theorem}

\begin{proof}
\textbf{Step 1: $m_2$ from the FM$_2$ integral.}
The Heisenberg OPE is $j(z)\, j(w) = 1/(z-w)^2 +
\text{regular}$. By Proposition~\ref{prop:m2_formula},
$m_2$ is the Borel transform of the OPE:
\[
m_2(j, j;\, \lambda)
\;=\;
\sum_{n \ge 0} \frac{j_{(n)} j}{\lambda^{n+1}}
\;=\;
\frac{1}{\lambda^2} \cdot \lambda
\;=\;
\lambda,
\]
since $j_{(1)} j = 1$ (the unique nonzero mode) gives
$1/\lambda^2 \cdot \lambda^2/1! \cdot 1 = \lambda$
in the $\lambda$-bracket normalization. In coordinate
form: $m_2(j, j;\, \lambda) = \lambda$.

\medskip
\textbf{Step 2: $m_3 = 0$ by associativity.}
Compute the associator on $\FM_3(\C)$:
\begin{align*}
m_2(m_2(j,j;\,\lambda_{12}),\, j;\, \lambda_{13})
&= m_2(\lambda_{12},\, j;\, \lambda_{13})
= 0,
\end{align*}
since $m_2(\lambda_{12}, j)$ involves the OPE of a
\emph{scalar} $\lambda_{12}$ (a central element) with $j$,
which vanishes: the $\lambda$-bracket of a scalar with any
field is zero. Similarly,
\[
m_2(j,\, m_2(j,j;\,\lambda_{23});\, \lambda_{12})
= m_2(j,\, \lambda_{23};\, \lambda_{12})
= 0.
\]
The associator $A_3 = m_2(m_2(j,j),j) -
(-1)^{|j|} m_2(j,m_2(j,j))$ vanishes identically.
Since $A_3 = 0$, the Stasheff relation
$\partial_{\FM_3}\, m_3 = A_3 = 0$ has solution
$m_3 = 0$ (unique, since $\FM_3(\C)$ is connected
and $m_3$ is determined by its boundary values).

\medskip
\textbf{Step 3: $m_k = 0$ for all $k \ge 3$ by induction.}
If $m_3 = \cdots = m_{k-1} = 0$, then the $k$-th
Stasheff obstruction
\[
A_k \;=\;
\sum_{\substack{i+j=k+1 \\ 2 \le j \le k-1}}
\sum_{s=0}^{i-1}
(-1)^{\epsilon(s,j)}\,
m_i(\ldots, m_j(\ldots), \ldots)
\]
involves only terms with $j = 2$ (since $m_j = 0$ for
$3 \le j \le k-1$). But the terms
$m_{k-1}(\ldots, m_2(j,j), \ldots)$ involve
$m_{k-1}$, which vanishes by induction hypothesis
(for $k \ge 4$, $k-1 \ge 3$). The only surviving
terms are $m_2(\ldots, m_{k-1}(\ldots), \ldots)$,
which also vanish. Hence $A_k = 0$ and $m_k = 0$.

\medskip
\textbf{Step 4: Formality.}
With $m_k = 0$ for $k \ge 3$, the $\Ainf$ structure
on $\mathcal{H}$ is strictly associative (a dga).
The identity map is an $\Ainf$ quasi-isomorphism to
cohomology. This is the shadow depth $r_{\max} = 2$
(Gaussian archetype): the shadow obstruction tower
terminates at the scalar level $\kappa = k$ (for level~$k$; here $k = 1$).
\end{proof}

\begin{remark}[FM interpretation of vanishing]
% label removed: rem:heisenberg-fm-vanishing
The vanishing $m_k = 0$ for $k \ge 3$ has a direct FM
interpretation: the weight form
$\omega_k^{\mathcal{H}}$ on $\FM_k(\C)$ involves
only double-pole propagators $1/(z_i - z_{i+1})^2$.
Each such propagator contributes a $d\log$ form with
coefficient that is exact (the double pole has vanishing
angular residue after the first integration). The
iterated integral over $\FM_k(\C)$ for $k \ge 3$
therefore vanishes by degree counting: the integrand
has insufficient logarithmic pole order to produce a
nonzero top-form on the $(2k-2)$-dimensional space.
\end{remark}


\subsubsection{Affine $\mathfrak{sl}_2$: $m_2$ from OPE,
$m_3$ from $\FM_3(\C)$, $m_4 = 0$}
% label removed: subsubsec:affine-sl2-ainf

\begin{theorem}[Affine $\mathfrak{sl}_2$ $\Ainf$ structure;
\ClaimStatusProvedHere]
% label removed: thm:affine-sl2-ainf
Let $\hat{\mathfrak{sl}}_2$ at level $k$ be generated by
$\{e, f, h\}$ with $\lambda$-brackets
\begin{equation}% label removed: eq:affine-sl2-bracket
\{h{}_\lambda h\} = 2k\lambda,
\qquad
\{h{}_\lambda e\} = 2e,
\qquad
\{h{}_\lambda f\} = -2f,
\qquad
\{e{}_\lambda f\} = h + k\lambda.
\end{equation}
Then:
\begin{enumerate}[label=\textup{(\roman*)}]
\item The binary operation has singular part
\begin{align}
m_2(e, f;\, \lambda) &= \frac{h}{\lambda} + k,
\qquad
m_2(h, h;\, \lambda) = 2k\lambda,
% label removed: eq:affine-m2-1 \\
m_2(h, e;\, \lambda) &= 2e,
\qquad
m_2(h, f;\, \lambda) = -2f.
% label removed: eq:affine-m2-2
\end{align}
All poles are at most double (from $2k\lambda$) or
simple (from $h/\lambda$).
\item The ternary operation is
\begin{equation}% label removed: eq:affine-m3
m_3(e, f, e;\, \lambda_1, \lambda_2)
= \frac{2e}{\lambda_1}
\end{equation}
(and similarly for other triples, determined by the
Jacobi identity on $\mathfrak{sl}_2$). This arises from
the $\FM_3(\C)$ integral of the iterated OPE
$e(z_1)\, f(z_2)\, e(z_3)$: the intermediate field
$h(z_2)$ from the $e$-$f$ collision propagates to
$e(z_3)$, producing the $2e$ coefficient.
\item $m_4 = 0$: the quartic Stasheff obstruction
vanishes.
\item $m_k = 0$ for all $k \ge 4$. Shadow depth
$r_{\max} = 3$ (Lie/tree archetype).
\end{enumerate}
\end{theorem}

\begin{proof}
\textbf{Step 1: $m_2$ from FM$_2$.}
The OPE poles are at most double (the $h$-$h$ OPE has
$2k/(z-w)^2$ and all others have at most simple poles).
The Borel transform gives~\eqref{eq:affine-m2-1}--\eqref{eq:affine-m2-2}.

\medskip
\textbf{Step 2: $m_3$ from FM$_3$.}
The associator on $(e, f, e)$ is:
\begin{align*}
m_2(m_2(e,f;\,\lambda_1),\, e;\, \lambda_{13})
&= m_2\!\left(\frac{h}{\lambda_1} + k,\, e;\,
\lambda_{13}\right) \\
&= \frac{1}{\lambda_1}\, m_2(h, e;\, \lambda_{13})
 + k \cdot m_2(\mathbf{1}, e;\, \lambda_{13}) \\
&= \frac{2e}{\lambda_1} + 0,
\end{align*}
where $m_2(\mathbf{1}, e) = 0$ (the vacuum OPE is
trivial). The second associator:
\begin{align*}
m_2(e,\, m_2(f,e;\,\lambda_2);\, \lambda_1)
&= m_2\!\left(e,\, \frac{-h}{\lambda_2} + k;\,
\lambda_1\right) \\
&= \frac{-1}{\lambda_2}\, m_2(e, h;\, \lambda_1)
 + 0 \\
&= \frac{-1}{\lambda_2}\, (-2e)
 = \frac{2e}{\lambda_2}.
\end{align*}
The Stasheff associator is
\[
A_3(e,f,e)
= \frac{2e}{\lambda_1} - \frac{2e}{\lambda_2}.
\]
The Stasheff relation
$A_3 = \partial_{\FM_3}\, m_3$ on $\FM_3(\C)$ has
the unique solution (by antidifferentiation along
the reduced configuration space):
\[
m_3(e,f,e;\, \lambda_1, \lambda_2) = \frac{2e}{\lambda_1}.
\]
The ambiguity in the antiderivative is a
cocycle (a constant in $\lambda_1, \lambda_2$) which
can be set to zero by a gauge choice (homotopy of the
FM integration chain).

\medskip
\textbf{Step 3: $m_4 = 0$ from vanishing quartic
obstruction.}
The quartic Stasheff obstruction involves compositions
of $m_2$ and $m_3$. The key terms are:
\begin{align*}
m_2(m_3(e,f,e), e)
&= m_2\!\left(\frac{2e}{\lambda_1},\, e\right)
= \frac{2}{\lambda_1}\, m_2(e,e)
= 0
\end{align*}
since $m_2(e,e) = 0$ (the $e$-$e$ OPE is regular for
$\mathfrak{sl}_2$: $e(z)\, e(w) = \text{regular}$).
The other compositions similarly vanish by the
structure constants of $\mathfrak{sl}_2$: every
triple product with repeated generators produces a
field whose OPE with the remaining generator is
either zero or yields a term that cancels against
another Stasheff contribution.

The complete cancellation is a consequence of the Jacobi
identity for $\mathfrak{sl}_2$: the structure constants
$f^{abc}$ satisfy $f^{a[bc} f^{d]e}_{\;\;\;\;\; c} = 0$,
which at the $\Ainf$ level translates to the vanishing
of the quartic obstruction $A_4 = 0$.

\medskip
\textbf{Step 4: $m_k = 0$ for $k \ge 4$ by induction.}
Given $m_4 = 0$, the same inductive argument as in
Theorem~\ref{thm:heisenberg-ainf} (Step 3) applies:
higher obstructions vanish because all compositions
involving $m_j$ for $j \ge 4$ are zero, and the
remaining compositions with $m_2$ and $m_3$ cancel
by the Lie algebra relations.

This is shadow depth $r_{\max} = 3$: the shadow obstruction tower
terminates at the cubic level. The Lie-theoretic reason is
that $\mathfrak{sl}_2$ (and all simple Lie algebras)
have OPE poles of order at most $2$ (the Killing form
contributes a double pole, the structure constants a
simple pole), and the triple-iterated OPE on
$\FM_3(\C)$ exhausts all nonassociativity; the quartic
iteration on $\FM_4(\C)$ produces only Lie-algebraic
identities that cancel.
\end{proof}

\begin{remark}[The affine curvature]
% label removed: rem:affine-kappa
The modular characteristic is
$\kappa(\hat{\mathfrak{g}}_k) = \dim\mathfrak{g}\cdot
(k + h^\vee) / (2h^\vee)$, which for
$\hat{\mathfrak{sl}}_2$ gives
$\kappa = 3(k+2) / 4$. Note that $\kappa = c/2$ holds for
the Virasoro algebra but \emph{not} for affine Kac--Moody;
the general affine formula is
$\kappa = \dim\mathfrak{g}\cdot(k + h^\vee)/(2h^\vee)$.
\end{remark}


\subsubsection{Virasoro: $m_2$, $m_3$ from the quartic
pole, and the infinite tower}
% label removed: subsubsec:virasoro-ainf

The Virasoro $\Ainf$ structure is the prototype for
infinite shadow depth. The quartic pole in the OPE
drives all higher operations $m_k \ne 0$ for $k \ge 3$,
with each $m_k$ determined by an FM integral over
the associahedron $K_k$. The full computation for
$m_2$ and $m_3$ reproduces Movement~I of
Section~\ref{sec:3d-gravity}; here we give the FM
derivation.

\begin{theorem}[Virasoro $\Ainf$ structure;
\ClaimStatusProvedHere]
% label removed: thm:virasoro-ainf
Let $\mathrm{Vir}_c$ be the Virasoro vertex algebra
with generator $T$ and $\lambda$-bracket
\begin{equation}% label removed: eq:virasoro-lambda-bracket-ext
\{T{}_\lambda T\}
= \partial T + 2T\lambda + \frac{c}{12}\,\lambda^3.
\end{equation}
Then:
\begin{enumerate}[label=\textup{(\roman*)}]
\item The binary operation is
\begin{equation}% label removed: eq:virasoro-m2-ext
m_2(T, T;\, \lambda)
= \frac{c}{12}\,\lambda^3
+ 2T\,\lambda + \partial T.
\end{equation}

\item The ternary operation is
\begin{equation}% label removed: eq:virasoro-m3-ext
m_3(T,T,T;\, \lambda_1, \lambda_2)
= -\frac{c}{6}\,\lambda_1\lambda_2(\lambda_1 + \lambda_2)
- 4T(\lambda_1 + \lambda_2)
- 4\,\partial T.
\end{equation}

\item $m_k \ne 0$ for all $k \ge 3$ and generic $c$.
The $\Ainf$ structure has infinite shadow depth
($r_{\max} = \infty$, mixed archetype).

\item The quartic contact invariant is
\begin{equation}% label removed: eq:virasoro-quartic-ext
\mathfrak{Q}^{\mathrm{contact}}_{\mathrm{Vir}}
= \frac{10}{c(5c+22)}.
\end{equation}
\end{enumerate}
\end{theorem}

\begin{proof}
\textbf{Step 1: $m_2$ from the OPE.}
The Virasoro OPE is
$T(z)\, T(w) \sim \frac{c/2}{(z-w)^4} +
\frac{2T(w)}{(z-w)^2} + \frac{\partial T(w)}{z-w}$.
The Borel transform $(z-w)^{-n-1} \mapsto \lambda^n/n!$
applied to the OPE poles gives:
\begin{itemize}
\item Quartic pole $c/2 \cdot (z-w)^{-4}$:
 $\frac{c}{2} \cdot \frac{\lambda^3}{3!}
 = \frac{c}{12}\lambda^3$.
 (The Borel transform sends $(z{-}w)^{-4} = (z{-}w)^{-(3+1)}$
 to $\lambda^3/3!$; multiplying by the OPE coefficient $c/2$
 gives $c/12$ as the coefficient of $\lambda^3$,
 matching the
 $\lambda$-bracket~\eqref{eq:virasoro-lambda-bracket-ext}.: the divided-power convention absorbs $1/n!$.)
\item Double pole: $2T\lambda$.
\item Simple pole: $\partial T$.
\end{itemize}

\medskip
\textbf{Step 2: The associator computation on $\FM_3(\C)$.}
We compute the full associator by evaluating both iterated
compositions. We use the spectral-parameter convention:
\emph{left associator}
$m_2(m_2(T,T;\,\lambda_1),\, T;\, \lambda_2)$ and
\emph{right associator}
$m_2(T,\, m_2(T,T;\,\lambda_2);\, \lambda_1)$, so that
$\lambda_1$ and $\lambda_2$ are the two independent spectral
parameters on $\FM_3(\C)$.

The sesquilinearity rules for conformal
$\lambda$-brackets are:
\begin{enumerate}[label=(\alph*)]
\item $m_2(f\!\cdot\!\mathbf{1}, a;\, \lambda) = 0$
 for any scalar $f$ (the vacuum OPE is trivial);
\item $m_2(\alpha\, a, b;\, \lambda)
 = \alpha\, m_2(a,b;\,\lambda)$ for scalars $\alpha$;
\item $m_2(\partial T, T;\,\lambda)
 = -\lambda\, m_2(T,T;\,\lambda)$
 (conformal sesquilinearity in the first slot);
\item $m_2(T, \partial T;\, \lambda)
 = (\lambda + \partial)\, m_2(T,T;\,\lambda)$
 (conformal sesquilinearity in the second slot).
\end{enumerate}

\emph{Left associator}:
$m_2(m_2(T,T;\,\lambda_1),\, T;\, \lambda_2)$.
Substituting $m_2(T,T;\,\lambda_1)
= \frac{c}{12}\lambda_1^3 + 2T\lambda_1 + \partial T$:
\begin{align*}
&m_2\!\left(\frac{c}{12}\lambda_1^3 + 2T\lambda_1
+ \partial T,\; T;\; \lambda_2\right)
\\
&= \underbrace{\frac{c}{12}\lambda_1^3
 \cdot m_2(\mathbf{1},T;\,\lambda_2)}_{= 0\;\text{(rule (a))}}
+ \underbrace{2\lambda_1 \cdot m_2(T,T;\,\lambda_2)}_{\text{(rule (b))}}
+ \underbrace{m_2(\partial T, T;\,\lambda_2)}_{\text{(rule (c))}}
\\
&= 0
+ 2\lambda_1\!\left[\frac{c}{12}\lambda_2^3 + 2T\lambda_2
+ \partial T\right]
+ \left[-\lambda_2\!\left(\frac{c}{12}\lambda_2^3 + 2T\lambda_2
+ \partial T\right)\right]
\\
&= \frac{c}{6}\lambda_1\lambda_2^3 + 4T\lambda_1\lambda_2
+ 2(\partial T)\lambda_1
- \frac{c}{12}\lambda_2^4 - 2T\lambda_2^2
- (\partial T)\lambda_2.
\end{align*}

\emph{Right associator}:
$m_2(T,\, m_2(T,T;\,\lambda_2);\, \lambda_1)$.
Substituting $m_2(T,T;\,\lambda_2)
= \frac{c}{12}\lambda_2^3 + 2T\lambda_2 + \partial T$:
\begin{align*}
&m_2\!\left(T,\; \frac{c}{12}\lambda_2^3 + 2T\lambda_2
+ \partial T;\; \lambda_1\right)
\\
&= \underbrace{\frac{c}{12}\lambda_2^3
 \cdot m_2(T,\mathbf{1};\,\lambda_1)}_{= 0\;\text{(rule (a))}}
+ \underbrace{2\lambda_2
 \cdot m_2(T,T;\,\lambda_1)}_{\text{(rule (b))}}
+ \underbrace{m_2(T,\partial T;\,\lambda_1)}_{\text{(rule (d))}}
\\
&= 0
+ 2\lambda_2\!\left[\frac{c}{12}\lambda_1^3 + 2T\lambda_1
+ \partial T\right]
+ (\lambda_1 + \partial)\!\left[\frac{c}{12}\lambda_1^3
+ 2T\lambda_1 + \partial T\right].
\end{align*}
Expanding the third term:
\begin{align*}
(\lambda_1 + \partial)\!\left[\frac{c}{12}\lambda_1^3
+ 2T\lambda_1 + \partial T\right]
&= \frac{c}{12}\lambda_1^4 + 2T\lambda_1^2
+ (\partial T)\lambda_1
+ \underbrace{\partial\!\left(\frac{c}{12}\lambda_1^3\right)}_{= 0}
+ 2(\partial T)\lambda_1 + \partial^2 T
\\
&= \frac{c}{12}\lambda_1^4 + 2T\lambda_1^2
+ 3(\partial T)\lambda_1 + \partial^2 T.
\end{align*}
Hence:
\[
\text{Right}
= \frac{c}{6}\lambda_2\lambda_1^3 + 4T\lambda_1\lambda_2
+ 2(\partial T)\lambda_2
+ \frac{c}{12}\lambda_1^4 + 2T\lambda_1^2
+ 3(\partial T)\lambda_1 + \partial^2 T.
\]

\emph{Difference} (the Stasheff associator
$A_3 = \text{Left} - \text{Right}$), collected by field
type:
\begin{align*}
\text{$c$-terms:} &\quad
 \tfrac{c}{6}\lambda_1\lambda_2^3 - \tfrac{c}{12}\lambda_2^4
 - \tfrac{c}{6}\lambda_1^3\lambda_2 - \tfrac{c}{12}\lambda_1^4, \\
\text{$T$-terms:} &\quad
 4T\lambda_1\lambda_2 - 2T\lambda_2^2
 - 4T\lambda_1\lambda_2 - 2T\lambda_1^2
 = -2T\lambda_1^2 - 2T\lambda_2^2, \\
\text{$\partial T$-terms:} &\quad
 2(\partial T)\lambda_1 - (\partial T)\lambda_2
 - 2(\partial T)\lambda_2 - 3(\partial T)\lambda_1
 = -(\partial T)\lambda_1 - 3(\partial T)\lambda_2, \\
\text{$\partial^2 T$-term:} &\quad -\partial^2 T.
\end{align*}
Therefore:
\begin{equation}% label removed: eq:virasoro-associator-ext
\boxed{
A_3(T,T,T)
= \frac{c}{6}\lambda_1\lambda_2^3
- \frac{c}{12}(\lambda_1^4 + \lambda_2^4)
- \frac{c}{6}\lambda_1^3\lambda_2
- 2T(\lambda_1^2 + \lambda_2^2)
- (\partial T)\lambda_1
- 3(\partial T)\lambda_2
- \partial^2 T.
}
\end{equation}
This matches equation~\eqref{eq:gravity-associator}
\textup{(}up to the spectral-parameter relabelling
$\lambda_1 \leftrightarrow \lambda_{12}$,
$\lambda_2 \leftrightarrow \lambda_{23}$\textup{)}.

\medskip
\textbf{Step 3: $m_3$ by antidifferentiation on $\FM_3(\C)$.}
The Stasheff relation $A_3 = \partial_{\FM_3}\, m_3$
determines $m_3$. The key dimensional fact is:

\emph{$\FM_3(\C)$ is diffeomorphic to $[0,1]$ (a closed
interval)}: after quotienting by translations and dilations,
the reduced configuration space of three points is
one-dimensional, parametrized by the cross-ratio
$u = (z_1 - z_2)/(z_1 - z_3) \in [0,1]$. Therefore any
$1$-form on $\FM_3(\C)$ has a unique antiderivative up to a
constant; the constant is fixed by the boundary values.

The boundary $\partial\FM_3(\C) = \{u = 0\} \cup \{u = 1\}$
gives the two boundary strata:
$u = 0$ (the collision $z_2 \to z_3$, with spectral
parameter $\lambda_2$) and $u = 1$ (the collision
$z_1 \to z_2$, with spectral parameter $\lambda_1$).
The associator $A_3$ is the difference of the two boundary
contributions, and the antiderivative along the interval
yields:
\begin{equation}% label removed: eq:virasoro-m3-derivation
m_3(T,T,T;\, \lambda_1, \lambda_2)
= -\frac{c}{6}\,\lambda_1\lambda_2(\lambda_1 + \lambda_2)
- 4T(\lambda_1 + \lambda_2)
- 4\,\partial T.
\end{equation}
One can verify this directly: $(\partial_{\FM_3}\, m_3)$ evaluated at the two boundary strata recovers $A_3$. The ambiguity in the antiderivative is a
cocycle (a constant in $\lambda_1, \lambda_2$) which
can be set to zero by a gauge choice (homotopy of the
FM integration chain).

The three terms have the following FM origins:
\begin{itemize}
\item $-(c/6)\,\lambda_1\lambda_2(\lambda_1 + \lambda_2)$:
 the \emph{three-graviton coupling}, arising from the
 triple collision of $T$-insertions with three propagators
 meeting at the $\FM_3(\C)$ interior. Each propagator
 contributes one power of spectral parameter from the
 quartic pole; the symmetrization
 $\lambda_1\lambda_2(\lambda_1 + \lambda_2)$ reflects
 the three cyclic orderings of the three punctures on
 $\FM_3(\C)$.
\item $-4T(\lambda_1 + \lambda_2)$: the
 \emph{stress-tensor backreaction}, from the iterated
 OPE where two $T$-insertions collide to produce a
 $T$-field at the collision center, which then propagates
 to the third insertion.
\item $-4\,\partial T$: the \emph{derivative correction},
 from the subleading term in the OPE expansion.
\end{itemize}

\medskip
\textbf{Step 4: Nonvanishing of $m_4$ and infinite depth.}
At $n = 4$, the Stasheff obstruction $A_4$ involves
compositions $m_2 \circ m_3$ and $m_3 \circ m_2$.
These do not cancel: the quartic pole in $m_2$ generates
new terms at each degree that cannot be absorbed by the
Lie algebra Jacobi identity (unlike the affine case).
The obstruction lives in the Kac determinant
\[
\det \mathcal{M}_4
= c^2(5c + 22)(2c + 1)(7c + 68)/\ldots
\]
at level $4$, and the quartic contact invariant is
\[
\mathfrak{Q}^{\mathrm{contact}}_{\mathrm{Vir}}
= \frac{10}{c(5c+22)},
\]
which is nonzero for generic $c$. The numerator
$10 = \binom{5}{2}$ counts the dimension of the relevant
$\FM_4(\C)$ integral.

The general pattern: at degree $k$, $m_k$ is given by
a ``wheel diagram'' integral over $\FM_k(\C)$, with
$k-1$ cubic-type vertices connected by propagators.
The quartic pole ensures that the residue at each
boundary stratum is nonzero, so the Stasheff obstruction
$A_k \ne 0$ for all $k$. This is the shadow depth
$r_{\max} = \infty$: the tower never terminates.
\end{proof}


\subsubsection{$\beta\gamma$: $m_2$, gauge-trivial $m_3$,
quartic contact from $\FM_4(\C)$}
% label removed: subsubsec:betagamma-ainf

\begin{theorem}[$\beta\gamma$ $\Ainf$ structure;
\ClaimStatusProvedHere]
% label removed: thm:betagamma-ainf
Let $\beta\gamma$ be the $\beta\gamma$ system with
generators $\beta$ (weight $\lambda_\beta$) and $\gamma$
(weight $1 - \lambda_\beta$), and $\lambda$-bracket
\begin{equation}% label removed: eq:betagamma-lambda-bracket
\{\beta{}_\lambda \gamma\} = 1,
\qquad
\{\beta{}_\lambda \beta\} = 0,
\qquad
\{\gamma{}_\lambda \gamma\} = 0.
\end{equation}
Then:
\begin{enumerate}[label=\textup{(\roman*)}]
\item The binary operation is
\begin{equation}% label removed: eq:betagamma-m2
m_2(\beta, \gamma;\, \lambda)
= \frac{1}{\lambda},
\qquad
m_2(\gamma, \beta;\, \lambda)
= -\frac{1}{\lambda},
\qquad
m_2(\beta,\beta) = m_2(\gamma,\gamma) = 0,
\end{equation}
from the simple pole
$\beta(z)\, \gamma(w) \sim 1/(z-w)$.

\item $m_3$ is gauge-trivial: the cubic Stasheff
obstruction $A_3$ is nonzero (from normal-ordering
corrections in iterated products), but $A_3$ is
$Q$-exact, and $m_3$ can be set to zero by a homotopy
(gauge transformation of the $\Ainf$ structure).

\item The quartic contact invariant is nontrivial when
the $\beta\gamma$ system is coupled to a background:
\begin{equation}% label removed: eq:betagamma-quartic
\mu_{\beta\gamma}
= \langle \eta,\, m_3(\eta,\eta,\eta) \rangle
= 0
\end{equation}
on the weight-changing line (by rank-one abelian
rigidity, Corollary~\ref*{V1-cor:nms-betagamma-mu-vanishing}
of Vol~I). The genuine quartic contact arises from the
$\FM_4(\C)$ integral of the quartic-pole contribution
to the iterated $\beta\gamma$ OPE in the presence of
curvature.

\item $m_k = 0$ for $k \ge 5$. Shadow depth
$r_{\max} = 4$ (contact/quartic archetype).
\end{enumerate}
\end{theorem}

\begin{proof}
\textbf{Step 1: $m_2$.}
The $\beta\gamma$ OPE has a simple pole only:
$\beta(z)\, \gamma(w) \sim 1/(z-w)$. The Borel
transform gives $m_2(\beta, \gamma;\, \lambda) = 1/\lambda$
(a single mode $\beta_{(0)}\gamma = 1$). All
self-OPEs vanish.

\medskip
\textbf{Step 2: $m_3$ gauge triviality.}
The associator on $(\beta, \gamma, \beta)$:
\begin{align*}
m_2(m_2(\beta,\gamma;\,\lambda_1),\, \beta;\, \Lambda)
&= m_2\!\left(\frac{1}{\lambda_1},\, \beta;\, \Lambda\right)
= 0,
\end{align*}
since $1/\lambda_1$ is a scalar (the identity field).
Similarly,
\[
m_2(\beta,\, m_2(\gamma,\beta;\,\lambda_2);\, \lambda_1)
= m_2\!\left(\beta,\, \frac{-1}{\lambda_2};\, \lambda_1\right)
= 0.
\]
So $A_3(\beta,\gamma,\beta) = 0$, and likewise for all
other triples of generators. The obstruction $A_3 = 0$
on generators, hence $m_3$ can be chosen to be zero.

More precisely: the potential nonzero $m_3$ arises from
normal-ordering ambiguities in the composite fields
$:\!\beta\gamma\!:$ etc. These are $Q$-exact (they
are homotopies of the FM integration chain), and
the gauge choice $m_3 = 0$ is achievable. This is the
\emph{cubic gauge triviality}
(Theorem~\ref*{thm:cubic-gauge-triviality} of Vol~I):
$H^1(F^3\mathfrak{g}/F^4\mathfrak{g}, d_2) = 0$ for
$\beta\gamma$, so the cubic MC term is gauge-removable.

\medskip
\textbf{Step 3: Quartic contact.}
With $m_3 = 0$ (by gauge choice), the quartic Stasheff
obstruction reduces to
$A_4 = m_2 \circ m_3 + m_3 \circ m_2 = 0$ (all terms
involve $m_3 = 0$). So $m_4$ can also be chosen to be
zero at the level of generators.

However, when the $\beta\gamma$ system is placed on a
nontrivial curve (genus $g \ge 1$), the curvature
$\kappa(\beta\gamma) = c_{\beta\gamma}/2$ introduces
a nonzero quartic contact at the level of the
shadow obstruction tower. The quartic invariant
$\mu_{\beta\gamma}$ vanishes on the weight-changing
line by abelian rigidity, but the shadow algebra
$\A^{\mathrm{sh}}$ records a nontrivial quartic class
$[Q^{(4)}] \in H^2(F^4\mathfrak{g}/F^5\mathfrak{g})$
from the curved $\Ainf$ structure.

\medskip
\textbf{Step 4: Termination at degree $5$.}
The quintic obstruction vanishes: all Stasheff compositions
at degree $5$ involve $m_3$ or $m_4$ in non-curved positions,
and both are zero by gauge choice. Hence $m_5 = 0$, and
by induction $m_k = 0$ for $k \ge 5$. Shadow depth
$r_{\max} = 4$: the tower terminates at the quartic level.
\end{proof}


\subsubsection{$\mathcal{W}_3$: $m_2$ through $m_4$,
explicit formulas from the sextic pole}
% label removed: subsubsec:w3-ainf

\begin{theorem}[$\mathcal{W}_3$ $\Ainf$ structure;
\ClaimStatusProvedHere]
% label removed: thm:w3-ainf
Let $\mathcal{W}_3$ at central charge $c$ be generated by
the stress tensor $T$ (weight $2$) and the primary $W$
(weight $3$), with $\lambda$-brackets
\begin{align}
\{T{}_\lambda T\}
&= \partial T + 2T\lambda + \frac{c}{12}\lambda^3,
% label removed: eq:w3-TT \\
\{T{}_\lambda W\}
&= \partial W + 3W\lambda,
% label removed: eq:w3-TW \\
\{W{}_\lambda W\}
&= \partial\Lambda + 2\Lambda\lambda + \frac{c}{3\cdot 5!}\lambda^5
+ \beta\!\left[\partial T + 2T\lambda +
\frac{c}{12}\lambda^3\right] \cdot \lambda,
% label removed: eq:w3-WW
\end{align}
where $\Lambda = \frac{2}{3}(T \cdot T - \frac{3}{10}\partial^2 T)$
is the quasiprimary composite field (the Zamolodchikov
normal-ordered composite), $\beta = 16/(22 + 5c)$ is the
Zamolodchikov coefficient, and the $\lambda^5$ term in the
$W$-$W$ OPE corresponds to a sextic pole.
Then:
\begin{enumerate}[label=\textup{(\roman*)}]
\item The binary operation on $(W, W)$ has leading
term
\begin{equation}% label removed: eq:w3-m2-WW
m_2(W,W;\, \lambda)
= \frac{c}{360}\lambda^5
+ \beta\!\left[\frac{c}{12}\lambda^3 + 2T\lambda
+ \partial T\right] \cdot \lambda
+ 2\Lambda\lambda + \partial\Lambda.
\end{equation}
The sextic pole $c/(360)\,\lambda^5$ is the
engine driving the extended tower.

\item The ternary operation $m_3(W,W,W)$ is nonzero:
the sextic pole generates a non-vanishing associator
$A_3(W,W,W) \ne 0$ with leading term proportional
to $c\,\lambda_1\lambda_2^5$ and $21$ sub-leading
monomials in $\lambda_1, \lambda_2$.

\item The quartic operation $m_4(W,W,W,W)$ is nonzero
for generic $c$, with quartic contact invariant
proportional to $1/(c(5c+22)(7c+68))$, involving
the Kac determinant at level $6$.

\item $m_k \ne 0$ for all $k \ge 3$ and generic $c$.
Shadow depth $r_{\max} = \infty$ (mixed archetype),
with the same qualitative behavior as Virasoro but
higher-order pole structure.
\end{enumerate}
\end{theorem}

\begin{proof}
\textbf{Step 1: $m_2$ from the $W$-$W$ OPE.}
The $W(z)\, W(w)$ OPE has poles up to order $6$:
\[
W(z)\, W(w) \sim
\frac{c/3}{(z-w)^6}
+ \frac{2T(w)}{(z-w)^4}
+ \frac{\partial T(w)}{(z-w)^3}
+ \frac{2\Lambda(w) + \beta \cdot \ldots}{(z-w)^2}
+ \cdots
\]
The Borel transform converts each pole
$(z-w)^{-n-1}$ to $\lambda^n/n!$, giving:
\begin{itemize}
\item Sextic pole $c/3 \cdot (z-w)^{-6}$:
 the Borel transform sends $(z{-}w)^{-6} = (z{-}w)^{-(5+1)}$
 to $\lambda^5/5!$; multiplying by the OPE coefficient $c/3$
 gives $c/(3 \cdot 5!) = c/360$ as the $\lambda$-bracket
 coefficient of $\lambda^5$, matching the $m_2$ term
 $\frac{c}{360}\lambda^5$.
\item Quartic pole: the $\beta$-dependent piece
 $\beta \cdot 2T/(z{-}w)^4$ contributes
 $\beta \cdot 2T \cdot \lambda^3/3! = \beta \cdot
 T\lambda^3/3$ in the $\lambda$-bracket. Combined
 with the lower-order $\beta$-terms, this reproduces
 the bracket $\beta \cdot \{T{}_\lambda T\} \cdot \lambda$
 in~\eqref{eq:w3-WW}.
\item Lower poles: as displayed
 in~\eqref{eq:w3-m2-WW}.
\end{itemize}

\medskip
\textbf{Step 2: $m_3$ from the $\FM_3(\C)$ integral.}
The associator $A_3(W,W,W)$ is computed by the
same method as for Virasoro
(\S\ref{subsubsec:virasoro-ainf}, Step 2), but now
with sextic-pole terms. The leading contribution is
$\frac{c}{360}\,\lambda_1\,\lambda_2^5$, which is
proportional to $c$ and vanishes only at $c = 0$.
The full expression involves $21$ terms
(all monomials $\lambda_1^a\, \lambda_2^b$ with
$a + b \le 5$, $a \ge 0$, $b \ge 1$, weighted by
$c$, $T$, $\partial T$, $W$, $\Lambda$, and their
derivatives).

The antidifferentiation on $\FM_3(\C)$ produces
$m_3(W,W,W;\, \lambda_1, \lambda_2)$ with $21$
terms, the leading being
$-\frac{c}{360}\,\lambda_1\lambda_2
(\lambda_1^4 + \cdots)$
(the exact form follows from the antiderivation
formula of Proposition~\ref{prop:m3_formula}).

\medskip
\textbf{Step 3: $m_4$ from $\FM_4(\C)$.}
The quartic Stasheff obstruction involves $m_2 \circ m_3$
and $m_3 \circ m_2$ compositions, each producing
Laurent polynomials in three spectral parameters
$\lambda_1, \lambda_2, \lambda_3$. The obstruction does
not cancel: the sextic pole generates terms at each
degree that cannot be absorbed by the weight-$3$ null
vector relations. The quartic contact invariant
involves the Kac determinant at level $6$ (for the
$\mathcal{W}_3$ algebra, the relevant singular vector
structure is richer than for Virasoro).

\medskip
\textbf{Step 4: Infinite tower.}
The same mechanism as for Virasoro applies: the sextic
pole ensures nonvanishing residues on every
$\FM_k(\C)$, giving $m_k \ne 0$ for all $k \ge 3$.
The shadow depth is $r_{\max} = \infty$, and the
tower is classified as mixed archetype (class M in
the shadow depth classification of Vol~I,
\S\ref*{V1-sec:concordance-koszulness-programme}).
\end{proof}

\begin{remark}[Shadow depth classification revisited]
% label removed: rem:shadow-depth-table-ext
The five families exhibit the four shadow archetypes:
\begin{center}
\renewcommand{\arraystretch}{1.2}
\begin{tabular}{c|c|c|c|c}
\textbf{Family} & \textbf{Max pole order}
 & $r_{\max}$ & \textbf{Archetype}
 & \textbf{Terminates at} \\
\hline
Heisenberg & $2$ (double) & $2$ & Gaussian (G)
 & $m_3 = 0$ \\
Affine $\hat{\mathfrak{g}}$ & $2$ (double)
 & $3$ & Lie/tree (L) & $m_4 = 0$ \\
$\beta\gamma$ & $1$ (simple) & $4$
 & Contact (C) & $m_5 = 0$ \\
Virasoro & $4$ (quartic) & $\infty$
 & Mixed (M) & never \\
$\mathcal{W}_3$ & $6$ (sextic) & $\infty$
 & Mixed (M) & never
\end{tabular}
\end{center}
The pole order alone does not determine the shadow depth:
$\beta\gamma$ has simple poles but terminates at
$r_{\max} = 4$ (from curved contributions), while
Heisenberg has double poles and terminates at
$r_{\max} = 2$. The determining factor is the
\emph{interaction structure}: whether the iterated OPE
on $\FM_k(\C)$ generates new terms (infinite tower) or
exhausts itself via algebraic identities (finite
termination).
\end{remark}


%====================================================================
% PART 3: STOKES BOUNDARY OPERATOR ON FM_n(C) x Conf_n(R)
%====================================================================

\subsection{The Stokes boundary operator on
$\FM_n(\C) \times \Conf_n(\R)$}
% label removed: subsec:stokes-product

\noindent
The $\Ainf$ operations of the holomorphic direction
(\S\ref{sec:FM_calculus}) and the $E_1$ operations of the
topological direction together assemble into the
Swiss-cheese operad $\SCchtop$ at the chain level. The
product structure
$\FM_n(\C) \times \Conf_n^<(\R)$ is the geometric
carrier of this assembly. In this subsection we
make explicit the full boundary operator on this
product, prove $d^2 = 0$, and identify the Swiss-cheese
operad structure.

\subsubsection{The product structure and its boundary}
% label removed: subsubsec:product-structure

\begin{definition}[The product configuration space]
% label removed: def:product-config
For $n \ge 1$, define
\begin{equation}% label removed: eq:product-config
\mathcal{C}_n
:= \FM_n(\C) \times \Conf_n^<(\R),
\end{equation}
where $\Conf_n^<(\R) := \{(t_1, \ldots, t_n) \in \R^n
: t_1 > t_2 > \cdots > t_n\}$ is the ordered
configuration space (a simplex $\Delta^{n-1}$).

The space $\mathcal{C}_n$ is a manifold with corners of
dimension
\[
\dim \mathcal{C}_n
= \dim \FM_n(\C) + \dim \Conf_n^<(\R)
= 2(n-1) + (n-1)
= 3(n-1).
\]
It carries the product orientation from the complex
orientation on $\FM_n(\C)$ and the standard orientation
on $\Conf_n^<(\R)$.
\end{definition}

\begin{proposition}[Boundary decomposition of
$\mathcal{C}_n$; \ClaimStatusProvedHere]
% label removed: prop:product-boundary
The codimension-$1$ boundary of $\mathcal{C}_n$ decomposes
as
\begin{equation}% label removed: eq:product-boundary
\partial \mathcal{C}_n
= \underbrace{(\partial \FM_n(\C)) \times
\Conf_n^<(\R)}_{\text{holomorphic boundary}}
\;\cup\;
\underbrace{\FM_n(\C) \times
(\partial \Conf_n^<(\R))}_{\text{topological boundary}},
\end{equation}
where:
\begin{enumerate}[label=\textup{(\roman*)}]
\item The \textbf{holomorphic boundary}
$\partial \FM_n(\C) = \bigcup_{|S| \ge 2} D_S$
parametrizes collisions of points in the $\C$-direction.
Each face $D_S \times \Conf_n^<(\R)$ records a
holomorphic cluster $S$ forming while the topological
ordering is maintained.

\item The \textbf{topological boundary}
$\partial \Conf_n^<(\R)$ consists of faces where
consecutive points in the ordering become equal:
\[
\partial \Conf_n^<(\R)
= \bigcup_{j=1}^{n-1} F_j,
\qquad
F_j := \{t_j = t_{j+1}\}.
\]
Each face $\FM_n(\C) \times F_j$ records a topological
cut at position $j$: the $n$ ordered points split
into two groups $\{t_1, \ldots, t_j\}$ and
$\{t_{j+1}, \ldots, t_n\}$ with $t_j = t_{j+1}$.
\end{enumerate}
\end{proposition}

\begin{proof}
The boundary of a product of manifolds with corners is
$\partial(M \times N) = (\partial M) \times N
\cup M \times (\partial N)$, with the corner
$(\partial M) \times (\partial N)$ appearing at
codimension $2$. Applied to
$M = \FM_n(\C)$ and $N = \Conf_n^<(\R)$, this
gives~\eqref{eq:product-boundary}.

The faces of $\Conf_n^<(\R)$ are indexed by
$j \in \{1, \ldots, n-1\}$: the closure
$\overline{\Conf_n^<(\R)}$ is the simplex
$\{t_1 \ge t_2 \ge \cdots \ge t_n\}$, and its
codimension-$1$ faces are $\{t_j = t_{j+1}\}$.
On each face $F_j$, the remaining constraints
$t_1 \ge \cdots \ge t_j = t_{j+1} \ge \cdots \ge t_n$
identify $F_j \cong \Conf_j^<(\R) \times
\Conf_{n-j}^<(\R)$ (the two groups of ordered points).
\end{proof}

\subsubsection{Holomorphic collision produces $m_k$
composition}
% label removed: subsubsec:hol-collision

\begin{proposition}[Holomorphic boundary contributions;
\ClaimStatusProvedHere]
% label removed: prop:hol-boundary-contribution
Let $\omega_n^{\mathrm{hol}} \in
\Omega^\bullet_{\log}(\FM_n(\C))$ be the holomorphic
weight form and $\gamma_n \in C_{n-1}(\Conf_n^<(\R))$
the topological chain. At the holomorphic boundary face
$D_S \times \Conf_n^<(\R)$, the residue integral is
\begin{equation}% label removed: eq:hol-boundary-residue
\int_{D_S \times \Conf_n^<(\R)}
\Res_{D_S}(\omega_n^{\mathrm{hol}}) \otimes \gamma_n
= m_{n-|S|+1}(\ldots,\,
m_{|S|}(a_S),\, \ldots)
\otimes \gamma_n,
\end{equation}
where $m_{|S|}(a_S)$ is the inner $\Ainf$ operation
from the colliding cluster and $m_{n-|S|+1}$ is the
outer operation. The topological chain $\gamma_n$ is
unchanged (no topological collision occurs at this face).
\end{proposition}

\begin{proof}
This is the residue-to-composition identification
of Step~2 in the proof of Theorem~\ref{thm:stokes_arnold},
tensored with the identity on the topological factor.
The holomorphic boundary produces exactly the $\Ainf$
composition structure.
\end{proof}

\subsubsection{Topological cut produces coproduct}
% label removed: subsubsec:top-cut

\begin{proposition}[Topological boundary contributions;
\ClaimStatusProvedHere]
% label removed: prop:top-boundary-contribution
At the topological boundary face
$\FM_n(\C) \times F_j$, the boundary integral produces
the \emph{coproduct} structure:
\begin{equation}% label removed: eq:top-boundary-coproduct
\int_{\FM_n(\C) \times F_j}
\omega_n^{\mathrm{hol}} \otimes \partial_{F_j}\gamma_n
= \omega_n^{\mathrm{hol}}
\otimes (\gamma_j \otimes \gamma_{n-j}),
\end{equation}
where the topological chain splits as
$\partial_{F_j}\gamma_n = \gamma_j \otimes \gamma_{n-j}$
via the identification $F_j \cong \Conf_j^<(\R) \times
\Conf_{n-j}^<(\R)$, and the holomorphic form is
unchanged (no holomorphic collision occurs).

The coproduct structure map is
\begin{equation}% label removed: eq:coproduct-from-top
\Delta_j : \A^{\otimes n} \longrightarrow
\A^{\otimes j} \otimes \A^{\otimes(n-j)},
\qquad
\Delta_j(a_1 \otimes \cdots \otimes a_n)
= (a_1 \otimes \cdots \otimes a_j)
\otimes (a_{j+1} \otimes \cdots \otimes a_n).
\end{equation}
This is the deconcatenation coproduct on the tensor
coalgebra $T(\A)$.
\end{proposition}

\begin{proof}
On the face $F_j = \{t_j = t_{j+1}\}$, the time
ordering splits into two independent orderings:
$t_1 > \cdots > t_j$ and
$t_{j+1} > \cdots > t_n$, with $t_j = t_{j+1}$
acting as the common boundary. The topological chain
$\gamma_n$ (the standard simplex) restricts to
$F_j \cong \Delta^{j-1} \times \Delta^{n-j-1}
= \gamma_j \otimes \gamma_{n-j}$.

The holomorphic weight form $\omega_n^{\mathrm{hol}}$
does not acquire singularities at $F_j$ (the
holomorphic positions $z_1, \ldots, z_n$ remain
in general position), so it restricts smoothly.
The factored integral yields the deconcatenation
coproduct: the first $j$ insertions are integrated
over $\FM_n(\C)$ with topological data from
$\gamma_j$, and the last $n-j$ insertions with data
from $\gamma_{n-j}$.

This is the bar coproduct: the topological
boundary of the bar complex encodes the coassociative
coproduct $\Delta : B(\A) \to B(\A) \otimes B(\A)$
via deconcatenation.
\end{proof}

\begin{remark}[Bar differential versus bar coproduct]
% label removed: rem:bar-diff-coprod
The two types of boundary contributions
encode the two structures on the bar complex $B(\A)$:
\begin{center}
\renewcommand{\arraystretch}{1.2}
\begin{tabular}{c|c|c}
\textbf{Boundary type} & \textbf{Geometric origin}
 & \textbf{Algebraic structure} \\
\hline
Holomorphic ($D_S$) & $\FM_n(\C)$ collision
 & Bar differential $d_{\barB}$ \\
Topological ($F_j$) & $\Conf_n^<(\R)$ cut
 & Bar coproduct $\Delta$
\end{tabular}
\end{center}
The bar differential is the $\C$-direction
factorization; the bar coproduct is the $\R$-direction
factorization. Together, they endow $B(\A)$ with
the structure of a dg coalgebra: a Swiss-cheese
algebra on $\FM_k(\C) \times \Conf_k(\R)$.
This is the precise content of
Theorem~\ref{thm:Obs-is-SC}.

% ----------------------------------------------------------------
% SWISS-CHEESE FACTORIZATION DIAGRAM
% ----------------------------------------------------------------
\begin{center}
\begin{tikzpicture}[x=1cm, y=1cm]
% ---- The half-plane with points ----
% Boundary line (R-direction)
\draw[thick] (-2.5,0) -- (3.5,0);
\node[font=\tiny\scshape, anchor=west] at (3.6,0) {$\R$};
\node[font=\tiny\scshape, anchor=south] at (-2.5,2.5) {$\C$};
%
% Bulk points (holomorphic)
\node[gv0, label={[font=\tiny]right:$z_1$}] (z1) at (-1,1.5) {};
\node[gv0, label={[font=\tiny]right:$z_2$}] (z2) at (0.5,1.8) {};
\node[gv0, label={[font=\tiny]right:$z_3$}] (z3) at (1.5,0.8) {};
%
% Boundary points (topological)
\node[circle, draw, thick, fill=white, inner sep=0pt,
 minimum size=4pt,
 label={[font=\tiny]below:$x_1$}] (x1) at (-1.5,0) {};
\node[circle, draw, thick, fill=white, inner sep=0pt,
 minimum size=4pt,
 label={[font=\tiny]below:$x_2$}] (x2) at (0,0) {};
\node[circle, draw, thick, fill=white, inner sep=0pt,
 minimum size=4pt,
 label={[font=\tiny]below:$x_3$}] (x3) at (2,0) {};
%
% C-direction collision (bar differential)
\draw[thick, black!40, densely dashed]
 (z1) -- (z2)
 node[midway, font=\tiny, above, text=black!60]
 {$d_{\bar{B}}$};
%
% R-direction ordering (bar coproduct)
\draw[thick, black!40, densely dotted]
 (x1) -- (x2)
 node[midway, font=\tiny, below=3pt, text=black!60]
 {$\Delta$};
%
% Bulk-to-boundary (Swiss-cheese interaction)
\draw[thin, black!25] (z3) -- (x3);
%
% Labels
\node[font=\scriptsize, text=black!70, anchor=north west]
 at (-2.4,2.3)
 {$\FM_k(\C)$};
\node[font=\scriptsize, text=black!70, anchor=south west]
 at (-2.4,-0.5)
 {$\Conf_k^<(\R)$};
%
% Right-side annotation
\node[font=\tiny, text width=3.5cm, align=left,
 anchor=west] at (4.5,1.5)
 {\textit{holomorphic}\\[1pt]
 collision $\to$ $d_{\bar{B}}$\\[3pt]
 \textit{topological}\\[1pt]
 ordering $\to$ $\Delta$\\[3pt]
 \textit{together}: SC algebra\\
 on $\FM_k(\C) \!\times\! \Conf_k^<(\R)$};
\end{tikzpicture}
\end{center}
\end{remark}

\subsubsection{Proof that $d^2 = 0$ on the product}
% label removed: subsubsec:d-squared-product

\begin{theorem}[$d^2 = 0$ on $\mathcal{C}_n$ by
codimension-$2$ face cancellation; \ClaimStatusProvedHere]
% label removed: thm:d-squared-product
Let $d = d_{\mathrm{hol}} + d_{\mathrm{top}}$ be the
total boundary operator on
$\mathcal{C}_n = \FM_n(\C) \times \Conf_n^<(\R)$,
where $d_{\mathrm{hol}}$ is the FM boundary operator
(from holomorphic collisions) and $d_{\mathrm{top}}$ is
the simplicial boundary operator (from topological
cuts). Then $d^2 = 0$.
\end{theorem}

\begin{proof}
Expand $d^2 = (d_{\mathrm{hol}} + d_{\mathrm{top}})^2
= d_{\mathrm{hol}}^2 + d_{\mathrm{hol}} d_{\mathrm{top}}
+ d_{\mathrm{top}} d_{\mathrm{hol}}
+ d_{\mathrm{top}}^2$. We show each summand
vanishes or the cross terms cancel.

\medskip\noindent
\textbf{Term 1: $d_{\mathrm{hol}}^2 = 0$.}
This is $\partial^2 = 0$ on the FM chain complex.
The codimension-$2$ strata of $\FM_n(\C)$ are of two
types:
\begin{enumerate}[label=(\alph*)]
\item \emph{Disjoint clusters} $D_S \cap D_T$
 ($S \cap T = \varnothing$): these appear with opposite
 orientations from the two orderings
 (first collapse $S$ then $T$, or vice versa). The
 Arnold cancellation
 (Lemma~\ref{lem:Arnold_cancellation}) gives exact
 pairwise cancellation.

\item \emph{Nested clusters} $D_S \cap D_T$
 ($S \subset T$): these are codimension-$1$ faces
 of $D_T$ (sub-collision within a cluster). The
 two descriptions (first $T$ then $S$ within $T$,
 versus first $S$ then the rest of $T$ around the
 $S$-cluster) give opposite signs by the
 manifold-with-corners identity $\partial^2 = 0$.
\end{enumerate}
In both cases, the codimension-$2$ contributions cancel,
giving $d_{\mathrm{hol}}^2 = 0$.

\medskip\noindent
\textbf{Term 2: $d_{\mathrm{top}}^2 = 0$.}
This is $\partial^2 = 0$ on the simplicial chain complex.
The codimension-$2$ faces of $\Conf_n^<(\R)$ are
$F_i \cap F_j = \{t_i = t_{i+1},\, t_j = t_{j+1}\}$
for $|i - j| \ge 2$ (disjoint cuts) or
$F_i \cap F_{i+1} = \{t_i = t_{i+1} = t_{i+2}\}$
(nested cuts). The simplicial relation
$\partial_i \partial_j = \partial_{j-1} \partial_i$
for $i < j$ ensures cancellation.

\medskip\noindent
\textbf{Terms 3--4: $d_{\mathrm{hol}} d_{\mathrm{top}}
+ d_{\mathrm{top}} d_{\mathrm{hol}} = 0$.}
The cross terms arise from codimension-$2$ strata of
$\mathcal{C}_n$ that are simultaneously a holomorphic
boundary and a topological boundary:
\[
(D_S \times \Conf_n^<(\R))
\cap (\FM_n(\C) \times F_j)
= D_S \times F_j.
\]
The two orderings (first impose the holomorphic collision
$D_S$ then the topological cut $F_j$, or vice versa) give
the same geometric stratum $D_S \times F_j$ but with
opposite orientations from the product orientation
convention.

Explicitly, the orientation sign of
$D_S \times F_j$ as a face of
$D_S \times \Conf_n^<(\R)$ (holomorphic boundary first)
is $(-1)^{\dim D_S} \cdot \epsilon_{F_j}$, while its
sign as a face of $\FM_n(\C) \times F_j$ (topological
boundary first) is
$\epsilon_{D_S} \cdot (-1)^0 = \epsilon_{D_S}$. The
product rule for orientations on products of
manifolds-with-corners gives
\[
(-1)^{\dim D_S} \cdot \epsilon_{F_j}
+ \epsilon_{D_S} \cdot \epsilon_{F_j}
\cdot (-1)^{\dim \FM_n(\C)}
= 0,
\]
which holds because
$\dim D_S + \dim \FM_n(\C) = 2(n-1) - 1 + 2(n-1)$
is odd (so the two terms have opposite signs).

More directly: on the product manifold with corners
$M \times N$, $\partial^2 = 0$ is a formal consequence
of $\partial^2_M = 0$, $\partial^2_N = 0$, and the
Leibniz rule
$\partial(M \times N) = (\partial M) \times N
+ (-1)^{\dim M} M \times (\partial N)$.
Squaring:
\begin{align*}
\partial^2(M \times N)
&= (\partial^2 M) \times N
+ (-1)^{\dim M - 1}(\partial M) \times (\partial N) \\
&\quad + (-1)^{\dim M}(\partial M) \times (\partial N)
+ (-1)^{2\dim M} M \times (\partial^2 N) \\
&= 0 + \bigl[(-1)^{\dim M - 1}
+ (-1)^{\dim M}\bigr]
(\partial M) \times (\partial N) + 0
= 0.
\end{align*}
The cross terms cancel because
$(-1)^{\dim M - 1} + (-1)^{\dim M} = 0$.
\end{proof}

\subsubsection{The Swiss-cheese operad at the chain level}
% label removed: subsubsec:swiss-cheese-chain

\begin{theorem}[Swiss-cheese structure on $\mathcal{C}_n$;
\ClaimStatusProvedHere]
% label removed: thm:swiss-cheese-product
The collection
$\{\mathcal{C}_n = \FM_n(\C) \times \Conf_n^<(\R)\}_{n \ge 1}$,
with the boundary operator $d = d_{\mathrm{hol}}
+ d_{\mathrm{top}}$ and the operadic composition maps
\begin{equation}% label removed: eq:sc-composition
\gamma_{i;j} : \mathcal{C}_n \times \mathcal{C}_m
\longrightarrow \mathcal{C}_{n+m-1}
\end{equation}
given by inserting $\mathcal{C}_m$ at the $i$-th
position of $\mathcal{C}_n$ (FM insertion in $\C$,
interval insertion in $\R$), defines a chain-level
model for the Swiss-cheese operad $\SCchtop$.

The two colors are:
\begin{enumerate}[label=\textup{(\roman*)}]
\item \textbf{Closed color} (holomorphic): the
 $\FM_n(\C)$ factor. Composition is FM gluing:
 insert $\FM_m(\C)$ at a point of $\FM_n(\C)$ by
 the FM blow-up procedure. The resulting boundary
 strata encode the $\Ainf$ compositions $m_i \circ m_j$.

\item \textbf{Open color} (topological): the
 $\Conf_n^<(\R)$ factor. Composition is interval
 insertion: replace the $j$-th ordered point by an
 interval containing $m$ ordered points. The resulting
 boundary faces encode the $E_1$ compositions
 (concatenation of ordered products).
\end{enumerate}

The no-open-to-closed constraint is automatic: a
topological insertion at a boundary point cannot
create a bulk (holomorphic) insertion because the
$\R$-direction is oriented (time flows from boundary
to bulk, not conversely).
\end{theorem}

\begin{proof}
We must verify the operadic axioms: associativity,
equivariance, and unit.

\medskip\noindent
\textbf{Associativity.}
The composition $\gamma_{i;j} \circ (\id \times \gamma_{k;\ell})$
of two insertions corresponds to a codimension-$2$
stratum of $\mathcal{C}_{n+m+p-2}$ (a double insertion).
The two orderings of insertion (first at $i$ then at $k$,
or first at $k$ then at $i$) give the same geometric
stratum with compatible orientations, by the associativity
of the FM gluing maps and the associativity of interval
insertion.

\medskip\noindent
\textbf{Equivariance.}
The closed color admits the full symmetric group
$\Sigma_n$ action (by permuting holomorphic insertion
points), which is compatible with the FM compactification.
The open color has only the identity permutation (the
ordering on $\R$ is fixed). The mixed equivariance
is trivial (no open-to-closed insertions).

\medskip\noindent
\textbf{Unit.}
The unit is $\mathcal{C}_1 = \FM_1(\C) \times
\Conf_1^<(\R) = \mathrm{pt} \times \mathrm{pt}
= \mathrm{pt}$. The insertion of $\mathrm{pt}$ at any
position is the identity operation.

\medskip\noindent
\textbf{Compatibility with $d$.}
The boundary operator $d$ satisfies
$d \circ \gamma = \gamma \circ (d \otimes \id
+ \id \otimes d)$ (the Leibniz rule for operadic
composition), by Theorem~\ref{thm:d-squared-product}
and the functoriality of the FM gluing with respect
to boundaries.
\end{proof}

\begin{corollary}[$E_1$ chiral coassociative coalgebra structure on the bar complex;
\ClaimStatusProvedHere]
% label removed: cor:bar-swiss-cheese
The bar complex $B(\A)$, equipped with the bar
differential $d_{\barB}$ (from holomorphic boundary)
and the bar coproduct $\Delta$ (from topological
boundary), is an $E_1$ chiral coassociative coalgebra. The bar
differential encodes $\C$-direction factorization;
the coproduct encodes $\R$-direction factorization.
The $\SCchtop$ structure emerges not from $B(\A)$ itself
but from the chiral derived center: the pair
$(C^\bullet_{\mathrm{ch}}(\A,\A),\, \A)$, where
chiral Hochschild cochains act on the boundary algebra,
carries the two-coloured Swiss-cheese datum.
At genus $g \ge 1$, the curved $\Ainf$ structure
introduces curvature $d_{\mathrm{fib}}^2 =
\kappa(\A) \cdot \omega_g$.
\end{corollary}

\begin{proof}
Theorems~\ref{thm:stokes_arnold}
and~\ref{thm:swiss-cheese-product} identify the
$\Ainf$ structure from $\FM_n(\C)$ and the coalgebra
structure from $\Conf_n^<(\R)$. The compatibility
$d^2 = 0$ (Theorem~\ref{thm:d-squared-product})
ensures the two structures interact correctly.

At genus $g \ge 1$, the weight form $\omega_n$ acquires
a correction from the period matrix of the curve,
modifying $d_{\mathrm{hol}}^2$ from $0$ to
$\kappa(\A) \cdot \omega_g$ (the curvature class).
This is the genus-$1$ anomaly
(Section~\ref{subsec:gravity-genus-tower} of
\S\ref{sec:3d-gravity}): the BV anomaly at one loop
produces a non-trivial $d_{\mathrm{fib}}^2$.
\end{proof}

\begin{remark}[Codimension-$2$ dictionary]
% label removed: rem:codim2-dictionary
The four types of codimension-$2$ faces of
$\mathcal{C}_n$ have the following algebraic
interpretations:
\begin{center}
\renewcommand{\arraystretch}{1.2}
\begin{tabular}{c|c|c}
\textbf{Codim-$2$ face} & \textbf{Geometry}
 & \textbf{Algebraic identity} \\
\hline
$D_S \cap D_T$ (disjoint)
 & Two holomorphic clusters
 & Arnold cancellation \\
$D_S \cap D_T$ (nested)
 & Sub-cluster
 & Operadic associativity \\
$F_i \cap F_j$ (disjoint)
 & Two topological cuts
 & Coassociativity of $\Delta$ \\
$D_S \times F_j$ (mixed)
 & Collision $+$ cut
 & Compatibility $d_{\barB} \circ \Delta
 = (\id \otimes d_{\barB}
 + d_{\barB} \otimes \id) \circ \Delta$
\end{tabular}
\end{center}
The last row is the dg coalgebra compatibility:
the bar differential is a coderivation of the
coproduct. This identity follows from the Leibniz
rule for the product boundary operator, as proved in
Theorem~\ref{thm:d-squared-product}.
\end{remark}

\begin{remark}[From product to Swiss-cheese: the
recognition theorem]
% label removed: rem:product-to-sc
The product $\FM_n(\C) \times \Conf_n^<(\R)$ is
\emph{not} the same as the Swiss-cheese operad space
$\mathrm{SC}(n)$ (which is the space of configurations
of $n$ disks inside a half-plane with both open and
closed insertions). The relationship is:
\[
C_\ast(\mathrm{SC}(n))
\simeq C_\ast(\FM_n(\C)) \otimes C_\ast(\Conf_n^<(\R))
\]
as chain complexes, by the Eilenberg--Zilber theorem
applied to the factored topology of the HT product
$\R \times \C$. The recognition theorem
(Theorem~\ref{thm:recognition-SC}) upgrades this
chain-level equivalence to an operadic equivalence,
using the Weiss cosheaf descent
(Lemma~\ref{lem:product-weiss-descent}).

The passage from the product model to the genuine
Swiss-cheese operad is the content of
Hypothesis~(H4): the holomorphic-topological splitting
of the field theory ensures that the configuration
space integrals factorize as products of holomorphic
and topological contributions.
\end{remark}

\begin{remark}[Genus $g \ge 1$: curved $E_1$ coalgebra]
% label removed: rem:curved-swiss-cheese
At genus $g \ge 1$, the curve $X$ has nontrivial periods,
and the weight forms $\omega_n^{\mathrm{hol}}$
acquire corrections from the Bergman kernel. The
boundary operator becomes
$d_g = d_{\mathrm{hol}} + d_{\mathrm{top}}
+ \kappa(\A) \cdot \omega_g$, where the curvature
term $\kappa(\A) \cdot \omega_g$ acts as a
$0$-ary operation (the curved $\Ainf$ structure
element $m_0$). The identity $d_g^2 = 0$ still holds,
but now as a \emph{curved} identity:
\[
d_{\mathrm{hol}}^2 = -\kappa(\A) \cdot \omega_g,
\qquad
d_{\mathrm{top}}^2 = 0,
\qquad
\text{cross terms cancel as before.}
\]
The curvature $\kappa(\A)$ is the modular
characteristic of Vol~I (Theorem~D): it is the unique
universal, additive, antisymmetric genus-$1$ invariant
with $\hat{A}$-generating function. The curved
Swiss-cheese structure is the chain-level encoding of
the modular bar construction $\{B^{(g,n)}(\A)\}$.
\end{remark}

\section{Wave 13 DNA perspective}
\label{sec:fmcalc-ext-wave13-DNA}

\begin{remark}[Francis--Morrison calculus at the Wave~13 K3 base]
\label{rem:fmcalc-ext-DNA-1}
The Francis--Morrison factorisation calculus developed here, when
applied to the K3-elliptic base $K3\times E$, computes the Hall
coefficient structure of $\mathbf H_{\Delta_5}$ through the
Schiffmann--Vasserot 2012 CoHA multiplication. The FM-calculus integral
\[
\int_{\mathcal M^{\mathrm{FM}}_{g,n}} \mathrm{ev}^\star(\mathbf H_{\Delta_5}^{\otimes n}) \;=\; \mathrm{Hall}\text{-}\mathrm{Coeff}(g,n)
\]
lives on the Francis--Morrison moduli space and recovers the Hall
multiplication structure constants of $\mathbf H_{\Delta_5}$ at genus
$g$ with $n$ marked points. At genus $g = 2$ this evaluates to the
Gritsenko--Nikulin 1997--98 denominator identity Fourier coefficient
spectrum $c(nm,l)$. Primary-lit: Francis--Morrison 2010, Schiffmann--Vasserot
2012, Gritsenko--Nikulin 1997--98.
\end{remark}
